\documentclass[11pt,a4paper]{article}

\usepackage[margin=25mm]{geometry}
\usepackage[T1]{fontenc}
\usepackage[utf8]{inputenc}
\usepackage{lmodern}
\usepackage{xcolor}
\usepackage{microtype}
\usepackage{booktabs}
\usepackage{longtable}
\usepackage{enumitem}
\usepackage{fancyhdr}
\usepackage{titlesec}
\usepackage{hyperref}
\usepackage{listings}
\usepackage{upquote}

\definecolor{ink}{HTML}{18324B}
\definecolor{accent}{HTML}{007C91}
\definecolor{codebg}{HTML}{F3F6F7}
\hypersetup{colorlinks=true,linkcolor=ink,urlcolor=accent,citecolor=ink,pdftitle={Building CLEO externally on Levante},pdfauthor={Anirudh Arora}}
\lstdefinestyle{shell}{backgroundcolor=\color{codebg},basicstyle=\ttfamily\small,breaklines=true,columns=fullflexible,frame=single,framerule=0.2pt,rulecolor=\color{ink!30},xleftmargin=2mm,xrightmargin=2mm,aboveskip=1.0em,belowskip=1.0em}
\lstset{style=shell}
\setlist[itemize]{leftmargin=*,itemsep=3pt,topsep=4pt}
\setlist[enumerate]{leftmargin=*,itemsep=4pt,topsep=4pt}
\pagestyle{fancy}
\fancyhf{}
\lhead{CLEO external build guide}
\rhead{Anirudh Arora}
\cfoot{\thepage}
\titleformat{\section}{\large\bfseries\color{ink}}{\thesection}{0.75em}{}
\titleformat{\subsection}{\normalsize\bfseries\color{ink}}{\thesubsection}{0.6em}{}

\title{\vspace{-1.5em}\textbf{Building CLEO externally on Levante}\\[0.4em]
\large A practical, pinned-source workflow from my 2-D condensation--collision experiment}
\author{Anirudh Arora}
\date{6 October 2026}

\begin{document}
\maketitle
\thispagestyle{fancy}

\begin{abstract}
This note records the way I built and ran the CLEO \texttt{constthermo2d} example for the fixed-grid 2-D condensation--collision experiment in the accompanying \texttt{CLEO-SDM-Convergence} branch. The central choice is simple: the experiment repository and the CLEO source tree are kept separate. I clone one exact upstream CLEO revision into a clean directory, verify that revision before every build or run, and retain the experiment-specific configuration, job scripts, receipts, and analysis in my own repository. This keeps a run inspectable without copying or modifying the upstream code base.
\end{abstract}

\tableofcontents
\newpage

\section{What this guide is for}

This is not intended as a replacement for the CLEO documentation. It is a concrete working pattern for a researcher who wants to use an upstream CLEO example without putting their experiment code inside a CLEO clone. The pattern was used for the 2-D \texttt{constthermo2d} condensation--collision experiment on Levante.

The relevant project branch is
\url{https://github.com/Arora-Anirudh/CLEO-SDM-Convergence/tree/share/constthermo2d-fixedgrid-experiment-2026-10-06}.
It contains the complete experiment-side record: configuration, batch scripts, input generation, audits, compact analysis products, figures, and the final report. The large raw Zarr trajectory stores are deliberately not tracked in Git.

\subsection{The main idea}

There are two repositories with different responsibilities:

\begin{longtable}{@{}p{0.27\linewidth}p{0.67\linewidth}@{}}
\toprule
Location & Responsibility \\
\midrule
\endhead
My experiment repository & Versioned YAML configurations, Slurm wrappers, Python input generators and analysis, figures, receipts, and scientific documentation. \\
Clean external CLEO clone & The unmodified upstream code at one recorded commit, its maintained build helper, and the compiled executable. \\
\bottomrule
\end{longtable}

The separation matters. It means that an experiment is not silently affected when upstream \texttt{main} advances, and it also means that the external CLEO source can be inspected or replaced without mixing those changes into the scientific record.

\section{The specific configuration recorded in this branch}

The example here is the CLEO \texttt{constthermo2d} executable. The experiment itself is not necessary to understand the build pattern, but it makes the directory names and scripts concrete.

\begin{itemize}
  \item Upstream source: CLEO v0.69.1 source at commit \texttt{9224585d784e9371926db9d0a23891a4b0241e24}.
  \item Executable target: \texttt{const2d} from \texttt{examples/constthermo2d}.
  \item Build mode: Kokkos host threads, with the number of threads set by the Slurm request.
  \item Experiment configuration: \texttt{config/constthermo2d\_fixedgrid\_baseline\_v1.yaml}.
  \item Domain: 1,500 m by 1,500 m in the x--z plane, 20 by 20 grid boxes, with a 20 m transverse extent.
  \item Simulation period: 0--7,200 s, with output every 120 s.
\end{itemize}

The exact source commit is a scientific input. I would not replace it by a tag name alone or by a moving branch name in a production workflow.

\section{Repository map: where the files live}

The following is the shortest route through the branch. The names are intentionally verbose because they form part of the provenance record.

\begin{longtable}{@{}p{0.42\linewidth}p{0.52\linewidth}@{}}
\toprule
Path & What it does \\
\midrule
\endhead
\texttt{config/constthermo2d\_fixedgrid\_baseline\_v1.yaml} & Baseline physical and numerical configuration. It is copied and then materialised for each resolution. \\
\texttt{scripts/generate\_constthermo2d\_seeded\_input\_v1.py} & Generates native grid, superdroplet, and thermodynamic input files from the external CLEO source. It also checks the initial superdroplet count and multiplicities. \\
\texttt{scripts/materialize\_constthermo2d\_resolution\_config\_v1.py} & Writes a fresh case configuration with a specified number of superdroplets per populated grid box and a specified thread count. \\
\texttt{scripts/run\_constthermo2d\_seeded\_case\_v1.py} & Runs a single fresh case, captures input and executable hashes, and writes a case receipt. \\
\texttt{scripts/audit\_constthermo2d\_case\_streaming\_v3.py} & Reads one stored output time at a time and checks the ragged trajectory structure, IDs, radii, and grid-box accounting. \\
\texttt{scripts/levante/build\_constthermo2d\_fixedgrid\_v4.sbatch} & The final compile-only Slurm job used here. It calls CLEO's maintained Levante build helper and builds only \texttt{const2d}. \\
\texttt{scripts/levante/run\_constthermo2d\_fixedgrid\_smoke\_v4.sbatch} & A one-case model and audit gate. Run this before an ensemble. \\
\texttt{scripts/levante/prepare\_constthermo2d\_fixedgrid\_production\_configs\_v1.sbatch} & Creates the shared, per-resolution YAML files before an array starts. This avoids a concurrent write race. \\
\texttt{scripts/levante/run\_constthermo2d\_fixedgrid\_production\_member\_v1.sbatch} & Runs one member in the fixed-20 production ensemble. \\
\texttt{scripts/levante/analyze\_constthermo2d\_fixedgrid\_production\_array\_v1.sbatch} & Streams accepted raw Zarr trajectories into small per-member summaries. \\
\texttt{docs/runs/constthermo2d\_fixedgrid\_*.md} & Run receipts, decisions, and what was checked at each stage. \\
\bottomrule
\end{longtable}

The generic project CMake integration is separate: \texttt{src/extern/cleo/CMakeLists.txt} and \texttt{docs/decisions/0001-external-cleo.md} show a \texttt{FetchContent} approach. That is useful when CLEO must be built as a dependency of a project-owned executable. It is \emph{not} the v0.69.1 build route used for the 2-D experiment below.

\section{Before building: make a clean external source checkout}

Choose a source location that is separate from both the experiment checkout and the build directory. On Levante, I kept this under project home storage because the source is small, long lived, and should not be repeatedly cloned into scratch.

\begin{lstlisting}[language=bash]
# Example paths. Adapt only the base directory and account name.
export PROJECT_HOME=/home/m/<account>/SDM
export CLEO_SOURCE=$PROJECT_HOME/CLEO-v0691-constthermo2d-9224585d
export CLEO_COMMIT=9224585d784e9371926db9d0a23891a4b0241e24

git clone https://github.com/yoctoyotta1024/CLEO.git "$CLEO_SOURCE"
git -C "$CLEO_SOURCE" checkout --detach "$CLEO_COMMIT"

# These two checks should both produce the expected, clean state.
git -C "$CLEO_SOURCE" rev-parse HEAD
git -C "$CLEO_SOURCE" status --porcelain
\end{lstlisting}

The expected hash from \texttt{rev-parse HEAD} is the complete value stored in \texttt{CLEO\_COMMIT}. The second command should print nothing. If it does print file names, do not build from that tree until the change is understood. In this workflow, any experimental change belongs in the experiment repository, not as an unrecorded edit to the source clone.

\subsection{External dependencies}

The final build wrapper expects a working Levante toolchain and the YAC/YAXT dependency prefix required by the CLEO build. In the recorded run, the prefix was supplied through \texttt{CLEO\_YACYAXTROOT}. The exact absolute paths in the committed batch script are specific to my DKRZ account and should be treated as examples, not copied verbatim.

Before starting a build, check that the static libraries exist in the prefix you intend to use:

\begin{lstlisting}[language=bash]
test -r "$CLEO_YACYAXTROOT/yac/lib/libyac.a"
test -r "$CLEO_YACYAXTROOT/yaxt/lib/libyaxt.a"
\end{lstlisting}

If these tests fail, resolve the CLEO/YAC/YAXT installation first. Changing the build script to bypass the checks only moves the failure later into CMake or linking.

\section{Build the external CLEO source}

\subsection{What the recorded build job does}

The final build script is
\texttt{scripts/levante/build\_constthermo2d\_fixedgrid\_v4.sbatch}. Before it compiles, it checks all of the following:

\begin{enumerate}
  \item the source directory exists;
  \item the expected YAC/YAXT libraries are present;
  \item the build directory does not already exist;
  \item the source checkout is exactly the pinned commit and has no local edits.
\end{enumerate}

It then exports CLEO's own build variables, sources the maintained upstream helper \texttt{scripts/levante/bash/build\_cleo.sh}, and asks CMake to build only the \texttt{const2d} target. The target-specific build is useful: it makes the build purpose clear and avoids compiling unrelated examples.

\subsection{Adapting the Slurm wrapper}

Copy the committed script rather than starting from an empty batch file, then edit only the site and user-specific values near the top:

\begin{lstlisting}[language=bash]
cp scripts/levante/build_constthermo2d_fixedgrid_v4.sbatch \
   scripts/levante/build_my_constthermo2d.sbatch

# In build_my_constthermo2d.sbatch, set these for your own account:
#SBATCH --account=<your-account>
readonly ROOT=/home/m/<account>/SDM/my_constthermo2d
readonly SOURCE=/home/m/<account>/SDM/CLEO-v0691-constthermo2d-9224585d
readonly BUILD=/home/m/<account>/SDM/cleo_builds/const2d-threads-gcc
readonly YACYAXTROOT=/home/m/<account>/SDM/cleo_dependencies/yacyaxt/gcc
\end{lstlisting}

Keep \texttt{SOURCE}, \texttt{BUILD}, and the experiment run root distinct. The script intentionally refuses to reuse \texttt{BUILD}; that avoids accidentally mixing CMake caches, compiler flags, or a different CLEO revision.

Submit the compile-only job from the experiment checkout:

\begin{lstlisting}[language=bash]
sbatch scripts/levante/build_my_constthermo2d.sbatch
squeue --me
\end{lstlisting}

The successful job writes a small receipt in its run root and ends with
\texttt{CONSTTHERMO2D\_COMPILE\_ONLY\_SUCCESS}. It also records a SHA-256 hash of the executable. I treat this receipt as the answer to ``what executable did this experiment use?''

\subsection{The key environment variables}

The relevant build settings are shown below. They are set in the final wrapper, not guessed at runtime.

\begin{lstlisting}[language=bash]
export CLEO_BUILDTYPE=threads
export CLEO_COMPILERNAME=gcc
export CLEO_ENABLEDEBUG=false
export CLEO_PATH2CLEO="$SOURCE"
export CLEO_PATH2BUILD="$BUILD"
export CLEO_YACYAXTROOT="$YACYAXTROOT"
export CLEO_BUILD_FLAGS="-DCLEO_COUPLED_DYNAMICS=fromfile \
  -DCLEO_DOMAIN=cartesian -DCLEO_NO_ROUGHPAPER=ON \
  -DCLEO_NO_PYBINDINGS=ON"

source "$SOURCE/scripts/levante/bash/build_cleo.sh"
cmake --build "$BUILD" --target const2d --parallel "$SLURM_CPUS_PER_TASK"
\end{lstlisting}

There is one small but important shell detail in the wrapper: it sources the site profile with shell ``nounset'' disabled. CLEO's Levante helper appends to optional environment variables, and a strict nounset shell can stop it before CMake begins. The wrapper records this explicitly rather than hiding it as an unexplained workaround.

\section{Generate inputs outside the CLEO source tree}

The build is only half of the reproducibility story. For each model member, the input files need their own clean case directory. The generator
\texttt{scripts/generate\_constthermo2d\_seeded\_input\_v1.py} writes grid, initial-superdroplet, thermodynamic, setup, and output paths into a copy of the baseline YAML. It imports \texttt{cleopy} and the upstream \texttt{examples.constthermo2d} input generator from the external source checkout.

The generator refuses to reuse a case directory and writes \texttt{input\_receipt.json}. It also verifies that the generated superdroplets have positive multiplicities and that the requested count appears in each populated grid box. For this experiment, 120 grid boxes are initially populated, so
\[
N_{\mathrm{initial}} = 120\,N_{\mathrm{cell}}.
\]

For example, an \texttt{N\_cell=8} case begins with 960 superdroplets. This is a configuration fact, not a convergence conclusion.

The native-input-only Slurm gate is
\texttt{scripts/levante/run\_constthermo2d\_fixedgrid\_input\_preflight\_v1.sbatch}. It is deliberately cheap: it generates and validates inputs for two small resolutions without running CLEO.

\begin{lstlisting}[language=bash]
sbatch scripts/levante/run_constthermo2d_fixedgrid_input_preflight_v1.sbatch
\end{lstlisting}

Only continue after the log contains
\texttt{CONSTTHERMO2D\_INPUT\_PREFLIGHT\_V1\_PASS} and the expected receipts exist.

\section{Smoke test the executable before an ensemble}

The correct next step is one full model run, not a large array. The recorded smoke wrapper is
\texttt{scripts/levante/run\_constthermo2d\_fixedgrid\_smoke\_v4.sbatch}. It checks the executable, external Python, external CLEO commit, clean source state, and all required package files before creating a fresh case directory.

\begin{lstlisting}[language=bash]
sbatch scripts/levante/run_constthermo2d_fixedgrid_smoke_v4.sbatch

# After the job completes, inspect both scheduler streams and the final marker.
tail -60 /scratch/m/m301324/SDM/constthermo2d_fixedgrid_baseline_v1/logs/smoke-v4-<jobid>.out
tail -60 /scratch/m/m301324/SDM/constthermo2d_fixedgrid_baseline_v1/logs/smoke-v4-<jobid>.err
\end{lstlisting}

For a successful recorded run, the output ends with
\texttt{CONSTTHERMO2D\_N8\_SMOKE\_SUCCESS}. The wrapper then creates plots, an input/case receipt, a trajectory-integrity JSON file, and a small SHA-256 list. A Slurm state of \texttt{COMPLETED} by itself is not enough; the completion marker, expected files, and audit output should agree.

\section{From a smoke test to a controlled ensemble}

The production workflow is deliberately split into stages rather than asking a single array job to create shared files on demand.

\begin{enumerate}
  \item \textbf{Materialise configurations once.} Submit \texttt{prepare\_constthermo2d\_fixedgrid\_production\_configs\_v1.sbatch}. It writes one immutable YAML per resolution, records their hashes, and avoids an array-wide race.
  \item \textbf{Run one independent member per task.} \texttt{run\_constthermo2d\_fixedgrid\_production\_member\_v1.sbatch} is parameterised with \texttt{NCELL} and \texttt{EXPECTED\_INITIAL}; its task index selects the ensemble member. The fixed-20 design has 20 deterministic but independent input seeds at each resolution.
  \item \textbf{Audit each member.} The member wrapper invokes the streaming trajectory audit. Failed or incomplete trajectories do not enter the aggregate analysis.
  \item \textbf{Aggregate only accepted trajectories.} \texttt{analyze\_constthermo2d\_fixedgrid\_production\_array\_v1.sbatch} runs the compact analysis one member at a time. The raw Zarr store remains authoritative on scratch; the repository keeps portable summaries, tables, and figures.
\end{enumerate}

Here is the form of a member submission. The array size and concurrency cap should be chosen to fit the project allocation and the cluster's current policy.

\begin{lstlisting}[language=bash]
# Example: 20 members for N_cell = 512, each initially containing 120 * 512 SDs.
sbatch --array=1-20%20 \
  --export=ALL,NCELL=512,EXPECTED_INITIAL=61440 \
  scripts/levante/run_constthermo2d_fixedgrid_production_member_v1.sbatch
\end{lstlisting}

The committed wrapper uses 16 host threads and requests 2 GiB for a production member. Those values were measured for this case and are not a universal CLEO prescription. A new configuration should first be profiled with a small run and then use the observed memory and wall time to set a request.

\section{Where the model output and logs belong}

I use home/project storage for code, configurations, packages, and receipts; I use scratch for raw Zarr output and scheduler logs. The important distinction is that raw output is regenerable, large, and normally short-lived, while the information required to identify and interpret it should survive in the repository.

For the recorded workflow the principal roots were:

\begin{lstlisting}[language=bash]
HOME_ROOT=/home/m/m301324/SDM/constthermo2d_fixedgrid_baseline_v1
SCRATCH_ROOT=/scratch/m/m301324/SDM/constthermo2d_fixedgrid_baseline_v1
\end{lstlisting}

The committed scripts name these locations exactly because they are historical records. For a new project, replace them consistently with your own paths. Do not direct multiple experiment versions to the same output directory. Every wrapper here is written to refuse reuse of a case directory for that reason.

\section{How analysis is connected to the run}

The production aggregation script, \texttt{scripts/analyze\_constthermo2d\_fixedgrid\_production\_v1.py}, reads an accepted Zarr trajectory and produces one compact \texttt{.npz} summary plus one JSON audit per member. The summary contains:

\begin{itemize}
  \item time series of superdroplet counts, droplet-number and radius moments, liquid mass, mass-weighted terminal speed, and mass-weighted radius quantiles;
  \item number and mass DSDs on fixed 250-, 500-, and 1000-bin logarithmic radius meshes;
  \item selected 2-D liquid-mass fields at 0, 60, and 120 minutes;
  \item agreement checks against CLEO's native observer diagnostics.
\end{itemize}

This is the useful compromise for a Git repository: a colleague can inspect the configuration and analysis logic, regenerate figures when raw data are available, and see the published compact products without downloading hundreds of large trajectory stores.

The curated final products are under \texttt{results/analysis/constthermo2d\_fixedgrid\_*}. The detailed scientific reading is in \texttt{docs/analysis/constthermo2d\_fixedgrid\_full\_diagnostics\_v1.md} and the report in \texttt{docs/reports/Two\_Dimensional\_CLEO\_Condensation\_Collision\_Convergence\_Report\_2026-09-28\_layout\_fixed.docx}.

\section{Alternative: project-level CMake FetchContent}

There is a second legitimate external-CLEO pattern in this repository. The root \texttt{CMakeLists.txt} includes \texttt{src/extern/cleo/CMakeLists.txt}, which uses CMake \texttt{FetchContent} to download a specific CLEO commit during configuration. This is appropriate when the project contains its own C++ executable that needs to link against CLEO.

The advantage is that a normal CMake configure step obtains the dependency automatically. The cost is that the source ends up under the build tree, and the build recipe must be maintained alongside the application. The 2-D experiment did not need a project-owned C++ target, so I preferred the clean-source approach above and CLEO's maintained Levante build helper.

Whichever route is chosen, the non-negotiable part is the pin: use a full commit SHA, record it in the experiment configuration or receipt, and do not use an unpinned upstream branch for a result that needs to be reproduced later.

\section{Practical checklist}

\begin{enumerate}
  \item Clone the experiment repository and identify the exact branch and experiment configuration.
  \item Clone upstream CLEO into a separate directory and check out the recorded full commit SHA.
  \item Confirm the CLEO source is clean and that YAC/YAXT are discoverable before compiling.
  \item Create a fresh build directory and run a compile-only job for the one target needed.
  \item Save the build receipt and executable hash.
  \item Run an input-only gate, then one full smoke test with an audit.
  \item Inspect the job output, error stream, final success marker, receipts, and expected output files.
  \item Only then submit an ensemble. Materialise shared configurations before the array begins.
  \item Keep raw Zarr on scratch, but commit configurations, scripts, receipts, compact summaries, figures, and interpretation.
\end{enumerate}

\section{A final note on scope}

The build workflow establishes software and numerical provenance. It does not, by itself, establish a scientific convergence result or turn an idealised in-domain diagnostic into rainfall. In this particular 2-D configuration, the horizontal direction is periodic and the boundary condition is null, so a spatial liquid-mass lobe or fast-drop diagnostic should not be described as surface precipitation. Keeping these interpretation boundaries in the documentation is just as important as recording the source commit.

\vfill
\noindent\textit{This guide was prepared from the versioned scripts and provenance records in the shared \texttt{constthermo2d} branch. For questions about the experiment-specific choices, please consult the cited configuration and run records before adapting them.}

\end{document}
